\section*{Euklidischer Algorithmus}

\subsection*{Beispiel: Euklidischer Algorithmus}
$a=48$ und $b=-30$
\begin{align*}
	48 &= (-1) \cdot (-30) + 18 \\
	-30 &= (-2) \cdot 18 + 6 \\
	18 &= 3 \cdot 6 + 0
\end{align*}
Als letzter Rest $\neq 0$ ist $\ggT(48,-30) = 6$.

\subsection*{Erweiterter euklidischer Algorithmus}

\subsubsection*{Lemma von B\'ezout}
$a,b\in\Z$, nicht beide 0.

Dann $\exists s,t\in\Z: \ggT(a,b)=sa + tb$.

\subsubsection*{Manuelle Methode}
\begin{enumerate}
	\item Berechne $\ggT(a,b)$
	\item Schrittw. rückw. einsetzen
\end{enumerate}

\subsection*{Anwendung}

\subsubsection*{Multiplikatives Inverse in $\Z_n$}
$0\neq z \in \Z_n$ hat Inv. $\iff \ggT(z,n)=1$

EEA liefert Inv. als $s$.

\subsubsection*{Lineare diophantische Gleichung}
\[c_0 = c_1\cdot x_1 + \ldots + c_n \cdot x_n\]
hat Lösung $\iff \ggT(c_1,\ldots,c_n) \mid c_0$.
